\documentclass[abstracton,parskip=half]{scrreprt}
\usepackage[english,ngerman]{babel}
\usepackage[utf8]{inputenc}
\usepackage[T1]{fontenc}
\usepackage{graphicx}
\usepackage{hyperref}
\usepackage{todonotes}
\usepackage{listings}
\usepackage{url}

\lstdefinelanguage{scala}{
  morekeywords={%
          abstract,case,catch,class,def,do,else,extends,%
          false,final,finally,for,forSome,if,implicit,import,lazy,%
          match,new,null,object,override,package,private,protected,%
          return,sealed,super,this,throw,trait,true,try,type,%
          val,var,while,with,yield},
  otherkeywords={=>,<-,<\%,<:,>:,\#,@},
  sensitive=true,
  morecomment=[l]{//},
  morecomment=[n]{/*}{*/},
  morestring=[b]",
  morestring=[b]',
  morestring=[b]"""
}[keywords,comments,strings]



\hypersetup{
   hidelinks=true,
   pdfauthor={Julian Andres Klode},
   pdftitle={Bachelor Thesis: Performance and Interfaces of Datatype-Generic Scala Programs},
   pdfsubject={Programming Languages},
}

\author{Julian Andres Klode\\Matriculation number: 2403668\\\\Supervisor: Yufei Cai}
% Let's abuse the subject field for now
\subject{\includegraphics[height=.13\textheight]{graphics/Logo_uni}\\Programming Languages and Software Technology Group\\Prof. Ostermann\\[1.5cm]}
\title{Performance and Interfaces of Datatype-Generic Scala Programs}
\subtitle{Bachelor Thesis}
\date{}


\begin{document}
    \pagenumbering{roman}
    \maketitle{}
    \begingroup
\let\titlepage\par
\let\endtitlepage
\let
\selectlanguage{ngerman}
\begin{abstract}
Generische Programmierung ist eine Vorgehensweise, die es ermöglicht, Funktionen
über die Struktur eines Datentypes zu definieren, anstelle von eines spezifischen
Datentypes. Dies erlaubt es einer Funktion auf beliebigen Datentypen zu operieren,
selbst auf solchen, die dem Programmier der Funktion noch nicht bekannt waren.

In dieser Bachelorarbeit werde ich Portierungen von drei Haskell Bibliotheken
(LIGD, EMGM, Uniplate) nach Scala vorstellen und ihre Funktionalität und Geschwindigkeit miteinander
sowie mit Shapeless vergleichen.

Ich werde außerdem zeigen, dass LIGD in Scala erweitert werden kann, sodass sie
beinahe so mächtig wie EMGM ist, indem ich das Problem der Erweiterbarkeit
von Funktionen löse, welches bisherige Implementierungen stark beeinträchtige.
\end{abstract}
\selectlanguage{english}
\begin{abstract}
Generic programming is a technique that allows functions to be defined over
the shape of an type, rather than a specific type; thus allowing one function
to be used on objects of arbitrary types, even those that were not known to
the programmer of the function.

In this thesis, I will present ports of three Haskell libraries (LIGD, EMGM, Uniplate) to Scala,
and compare their features and performance against each other and
against shapeless, a native Scala library.

I will show that LIGD in Scala can be extended to be almost as powerful as EMGM;
solving the issue with extensibility that made it vastly inferior in previous
implementations.


\end{abstract}
\endgroup

    \selectlanguage{english}

    \tableofcontents{}
    \clearpage
    \pagenumbering{arabic}
    \setcounter{page}{1}

    \chapter{Introduction}

Many\footnote{Such as Haskell, Java, Scala, C++} programming languages these days support a form of parametric polymorphism.
Parametric Polymorphism allows the definition of data types that abstract over
the type of some element they store, such as lists that can store elements of
some unspecified type \lstinline{a}:
\begin{lstlisting}[language=scala]
sealed abstract class List[a]
case class Nil[a]() extends List[a]
case class Cons[a](car: a, cdr: List[a]) extends List[a]
\end{lstlisting}
In Scala and other object-oriented (or influenced) languages this concept is
called generics.
Datatype-generic programming is a different concept. Instead of abstracting
over the type of the elements, it abstracts over the container type.

For example, consider we have types to represent arrays and lists. Consider that
those store integers. We want to increase every integer by one. That is, we want
to write a function that takes a container of integers and increases every integer
by one. That is, we want to write something like:

\begin{lstlisting}[language=Scala]
def incInt[C[_]](c: C[Int]) = everywhere((i: Int) => i + 1)(c)
\end{lstlisting}
(Note: This does not actually work with shapeless)

Datatype-generic programming does not only abstract containers. In Haskell and
other functional programming languages, data types (except some built-in scalar
types like integers and floating point numbers) are represented using
algebraic data types, that is (labelled) sum types and product types.
In Scala, case classes with a sealed abstract base class (or trait) can
represent algebraic data types.\cite{scalagp}
But Scala also allows unsealed classes, so further cases can be added even in
different modules.


\section{Existing Haskell Modules}
\subsection{\texttt{Data.Generics}}
\subsection{\texttt{Unbound}}

\cite{DBLP:conf/tldi/LammelJ03}


    \chapter{Implementing Datatype-Generic Programming in Scala}
    \section{Name-binding Operations}
    \section{Derived Data Types}
    \section{GPBench}
    GPBench is a suite of benchmarks for comparing generic programming
    libraries in Haskell \cite{DBLP:conf/haskell/RodriguezJJGKO08}.

    \chapter{Conclusion}

    \bibliographystyle{alpha}
    \bibliography{thesis}
\end{document}
